\documentclass[11pt]{article}
\usepackage[margin=1in]{geometry}
\usepackage{amsmath,amssymb,amsthm}
\usepackage[T1]{fontenc}
\usepackage{lmodern}
\usepackage{xcolor}
\usepackage{booktabs}
\usepackage{array}
\usepackage{longtable}
\usepackage[hidelinks]{hyperref}
\usepackage{url}
\usepackage{newunicodechar}
\newunicodechar{∠}{\ensuremath{\angle}}
\newunicodechar{₁}{\ensuremath{_1}}
\newunicodechar{₂}{\ensuremath{_2}}
\newunicodechar{Ω}{\ensuremath{\Omega}}
\newunicodechar{∙}{\ensuremath{\cdot}}

\newtheorem{theorem}{Theorem}
\newtheorem{lemma}[theorem]{Lemma}
\theoremstyle{definition}
\newtheorem{definition}[theorem]{Definition}
\theoremstyle{remark}
\newtheorem{remark}[theorem]{Remark}

\newcommand{\lean}[1]{\texttt{#1}}
% A declaration of one of our developments; every \decl{...} must be in that repository's gate.sh
% REQUIRED list (checked by the name gate).
\DeclareUrlCommand\decl{\urlstyle{tt}\def\UrlBreaks{\do\_\do\.}}
% A Mathlib declaration (pinned revision 0df444a3); not one of ours.
\DeclareUrlCommand\mdecl{\urlstyle{tt}\def\UrlBreaks{\do\_\do\.}}
\newcommand{\RR}{\mathbb{R}}
\newcommand{\PH}[1]{\textcolor{red}{\textbf{[PLACEHOLDER:} #1\textbf{]}}}

% ---------------------------------------------------------------------------------------------
% PINS. Filled in when the repositories are published; until then \src prints file:line only.
% No commit is pinned yet: define \csname commitof:paperg_nagel\endcsname etc. at publication.
% ---------------------------------------------------------------------------------------------
\expandafter\def\csname commitof:paperg_nagel\endcsname{e40ee54f20588f5181d348767f3f31be20406cc1}
\expandafter\def\csname commitof:paperg_brocard\endcsname{ebb35755a81aec2c7c0673fe93baf8cc5d6eb22d}
\newcommand{\commitof}[1]{\csname commitof:#1\endcsname}
\newcommand{\repo}[1]{https://github.com/jpbald93/#1}
\newcommand{\src}[4]{\ifcsname commitof:#1\endcsname
  \href{\repo{#1}/blob/\commitof{#1}/#2/#3.lean\#L#4}{\texttt{#3}:#4}\else\texttt{#3}:#4\fi}
\title{Triangle centres in Lean~4:\\ the Nagel line family and the symmedian and Brocard points}
\author{Josh Bald\\ \small Independent researcher, Ontario, Canada\\ \small\texttt{jpbald93@gmail.com}\\
\small ORCID: \texttt{0009-0002-1317-6489}
\thanks{\copyright~2026 Josh Bald. Licensed under CC BY 4.0 (\url{https://creativecommons.org/licenses/by/4.0/}). The accompanying Lean code is licensed under the Apache License 2.0.}}
\date{\PH{date}}
\emergencystretch=3em
\begin{document}
\sloppy
\maketitle
\begin{abstract}
We give machine-checked proofs, in Lean~4 with Mathlib, of classical theorems about two families of
triangle centres. For the first family we define the Nagel and Gergonne points by their barycentric
weights and prove that they are the common points of the cevians to the excircle and incircle contact
points, where the contact points are Mathlib's \texttt{touchpoint}s. We prove the Nagel line
relation $N-G=2(G-I)$ with Mathlib's incentre and centroid, that the Spieker centre (the incentre of
Mathlib's medial triangle) is the midpoint of $I$ and $N$, that the Nagel point of the medial triangle
is $I$, that the incircle and excircle contact points on each side are symmetric about its midpoint,
that $N$ is the incentre of the anticomplementary triangle, that the de Longchamps point lies on the
Euler line and is the orthocentre of the anticomplementary triangle, and that the Spieker centre has
equal power with respect to the three excircles. For the second family we prove that the point with
barycentrics $a^2:b^2:c^2$ lies on the three symmedians (six angle equalities) and inside the triangle,
that it is the centroid of its pedal triangle and, in the plane, the only such point, that the two
Brocard points satisfy their angle equalities, that the Brocard angle $\omega$ satisfies
$\cot\omega=\cot A+\cot B+\cot C=(a^2+b^2+c^2)/(4\Delta)$ and $0<\omega\le\pi/6$, that the Brocard
points are isogonal and equidistant from the circumcentre, that $O$, $K$ and both Brocard points lie
on the circle with diameter $OK$ (which degenerates exactly for equilateral triangles), that the six
points of Lemoine's first circle are concyclic, and the Droz-Farny line theorem. Apart from the
pedal-centroid converse and Droz-Farny, no numbered theorem assumes that the ambient space is a plane. All proofs
use only Lean's three standard axioms. In the sources we searched, subject to stated gaps, we found
no earlier formal proof of these statements; one earlier Lean file proves an algebraic form of the
Nagel and Gergonne cevian statements, and an earlier Coq development verifies barycentric forms of
the Nagel line, the Spieker centre's position and two related facts.
\end{abstract}

\section{Introduction}\label{sec:intro}
Triangle centres are among the most studied objects of elementary geometry; Kimberling's
\emph{Encyclopedia of Triangle Centers} \cite{etc} lists tens of thousands of them. A handful are
classical, with theorems in every textbook \cite{johnson1929,coxetergreitzer,honsberger}: the
Gergonne point $X(7)$ and the Nagel point $X(8)$, where the cevians to the incircle and excircle
contact points meet; the Spieker centre $X(10)$; the de Longchamps point $X(20)$; the symmedian
(Lemoine) point $X(6)$; and the two Brocard points. This paper formalises the basic theorems about
these points in Lean~4 \cite{lean4} with Mathlib \cite{mathlib}.

Mathlib already contains a good part of the surrounding theory, and we build on it rather than
redoing it. The Euler line is there: the Monge point of a simplex is
$O+\tfrac{n+1}{n-1}(G-O)$ (\mdecl{Affine.Simplex.mongePoint_eq_smul_vsub_vadd_circumcenter}),
and for a triangle the orthocentre is $O+3(G-O)$
(\mdecl{Affine.Triangle.orthocenter_eq_smul_vsub_vadd_circumcenter}). Mathlib also has the
incentre and excentres of a simplex, their radii and spheres, and the contact points
\mdecl{Affine.Simplex.touchpoint} (the orthogonal projections of the in- or excentre onto the faces),
together with the medial simplex \mdecl{Affine.Simplex.medial}. Our definitions of the Nagel and
Gergonne points are by weights, but every theorem about them is stated with Mathlib's
\lean{incenter}, \lean{excenter}, \lean{touchpoint}, \lean{medial}, \lean{circumcenter} and
\lean{orthocenter}, so that the bridges between the weights and Mathlib's notions are part of what is
proved. Euler's relation $OI^2=R(R-2r)$ is the subject of an open Mathlib pull request
(\#44230 \cite{pr44230}, by another author); it is related work, not part of this paper, and we do not use it.

\subsection{Contributions}
The numbered theorems are proved for an arbitrary \mdecl{Affine.Triangle} in a real inner product space of any
dimension, with no hypothesis beyond the triangle itself, except where stated.
\begin{itemize}
\item \emph{Nagel line family} (Section~\ref{sec:nagel}). Gergonne and Nagel concurrence through
Mathlib's contact points (Theorems~\ref{thm:gergonne} and~\ref{thm:nagel}); the Nagel line
$N-G=2(G-I)$ (Theorem~\ref{thm:nagelline}); the Spieker centre as the midpoint of $IN$
(Theorem~\ref{thm:spieker}); the de Longchamps point (Theorem~\ref{thm:delong}); the medial and
anticomplementary triangles (Theorems~\ref{thm:medialnagel} and~\ref{thm:anti}); the isotomic
contact points (Theorem~\ref{thm:isotomic}); and the Spieker centre as the radical centre of the
excircles, with the value of the common power (Theorem~\ref{thm:radical}).
\item \emph{Symmedian and Brocard points} (Section~\ref{sec:brocard}). The six symmedian angle
equalities and interiority (Theorem~\ref{thm:symmedian}); the symmedian point as the centroid of its
pedal triangle, and its uniqueness in the plane (Theorem~\ref{thm:pedal}); the two Brocard points
(Theorem~\ref{thm:brocard}); the Brocard angle (Theorem~\ref{thm:cot}); isogonality and
equidistance from $O$ (Theorem~\ref{thm:isogonal}); the Brocard circle and its degenerate case
(Theorems~\ref{thm:brocardcircle} and~\ref{thm:degenerate}); Lemoine's first circle
(Theorem~\ref{thm:lemoine}); and the Droz-Farny line (Theorem~\ref{thm:droz}).
\end{itemize}
The proofs are by coordinate certificates, as in our earlier papers
\cite{paperA,paperB,paperC,paperD,paperE,paperF}: barycentric coordinates relative to the triangle,
the Ravi substitution for the Nagel family, and inner products in the basis $B-A$, $C-A$ for the
angle statements (Section~\ref{sec:method}). Table~\ref{tab:thms} lists the Lean name of every
numbered theorem.

\section{Setting and definitions}\label{sec:defs}
Throughout, $V$ is a real inner product space, $P$ a metric space that is a normed torsor over $V$,
and $t$ an \mdecl{Affine.Triangle} in $P$, that is, three affinely independent points
$A=t_0$, $B=t_1$, $C=t_2$. The dimension of $V$ is arbitrary unless stated otherwise. We write
\[
a=|BC|,\quad b=|CA|,\quad c=|AB|,\quad s=\tfrac12(a+b+c),
\]
$G$ for the centroid (Mathlib's \mdecl{Finset.centroid} of the vertices), $I$ for
\mdecl{Affine.Simplex.incenter}, $O$ for \lean{circumcenter} and $H$ for \lean{orthocenter}. For
weights $w_0,w_1,w_2$ with $w_0+w_1+w_2=1$ we write $w_0A+w_1B+w_2C$ for Mathlib's
\mdecl{Finset.affineCombination}; all the points below are defined this way, by normalised
barycentric coordinates. Because $t$ is non-degenerate, $s-a$, $s-b$, $s-c$ and $a,b,c$ are positive
(the library proves $s-a>0$ etc.\ from Mathlib's positivity of the heights, through Heron's
formula), so none of the normalisations divides by zero.

\paragraph{Mathlib's contact points.} For a set $S$ of vertex indices, Mathlib's
\lean{t.touchpoint S i} is the orthogonal projection of the (ex)centre with sign pattern $S$ onto the
face opposite vertex~$i$. With $S=\emptyset$ this is the incentre, and \lean{t.touchpoint \(\emptyset\) i}
is the contact point of the incircle with the side opposite vertex~$i$. We prove that
\lean{t.excenter \{i\}} is the excentre opposite vertex $i$ (its weights are
$(-a,b,c)/(2(s-a))$ for $i=0$, and cyclically), so \lean{t.touchpoint \{i\} i} is the contact point of
the excircle opposite vertex $i$ with the side opposite vertex $i$.

\begin{definition}[Nagel family]\label{def:nagel}
\begin{enumerate}
\item The \emph{Nagel point} is $N=\dfrac{(s-a)A+(s-b)B+(s-c)C}{s}$
(\decl{Affine.Triangle.nagelPoint}).
\item The \emph{Gergonne point} is the point with weights proportional to
$\dfrac1{s-a},\dfrac1{s-b},\dfrac1{s-c}$, normalised to sum~$1$
(\decl{Affine.Triangle.gergonnePoint}).
\item The \emph{Spieker centre} is the incentre of Mathlib's medial triangle
\mdecl{Affine.Simplex.medial} (\decl{Affine.Triangle.spiekerCenter}).
\item The \emph{de Longchamps point} is $L=H+2(O-H)$, the reflection of $H$ in $O$
(\decl{Affine.Triangle.deLongchampsPoint}).
\item The \emph{anticomplementary triangle} is the image of $t$ under the homothety with centre $G$
and ratio $-2$ (\decl{Affine.Triangle.anticomplementary}). It is defined as a value of type
\mdecl{Affine.Triangle} (via \lean{Simplex.map} and the injectivity of the homothety), so its
non-degeneracy is part of the construction; its vertices are $B+C-A$, $C+A-B$, $A+B-C$
(\decl{NagelLine.anticomplementary_points_zero} and its two companions), and its medial
triangle is $t$ (\decl{NagelLine.anticomplementary_medial}).
\end{enumerate}
\end{definition}

\begin{definition}[symmedian and Brocard points]\label{def:brocard}
\begin{enumerate}
\item The \emph{symmedian point} is $K=\dfrac{a^2A+b^2B+c^2C}{a^2+b^2+c^2}$
(\decl{Affine.Triangle.symmedianPoint}).
\item The \emph{first Brocard point} $\Omega$ has weights proportional to
$\dfrac1{b^2},\dfrac1{c^2},\dfrac1{a^2}$ on $A,B,C$, normalised by
$\dfrac1{a^2}+\dfrac1{b^2}+\dfrac1{c^2}$ (\texttt{Affine.Triangle.brocardPoint\textsubscript{1}}).
\item The \emph{second Brocard point} $\Omega'$ has weights proportional to
$\dfrac1{c^2},\dfrac1{a^2},\dfrac1{b^2}$, with the same normalisation
(\texttt{Affine.Triangle.brocardPoint\textsubscript{2}}).
\item The \emph{area} is $\Delta=\tfrac12\,|AB|\,|AC|\sin\angle BAC$ (\decl{Affine.Triangle.area}),
with Mathlib's unoriented angle \mdecl{EuclideanGeometry.angle}. The library also proves the Gram
form $\Delta=\tfrac12\sqrt{|u|^2|v|^2-\langle u,v\rangle^2}$ with $u=B-A$, $v=C-A$
(\decl{Brocard.area_eq_sqrt}) and Heron's form in $a^2,b^2,c^2$ (\decl{Brocard.area_eq_sqrt_dist}).
\item The six \emph{Lemoine points} (\decl{Affine.Triangle.lemoinePoints}): with
$S=a^2+b^2+c^2$, they are the points of the side lines with parameters
$C+\tfrac{a^2}{S}(A-C)$, $A+\tfrac{b^2+c^2}{S}(B-A)$, $B+\tfrac{c^2+a^2}{S}(C-B)$,
$A+\tfrac{b^2}{S}(B-A)$, $B+\tfrac{c^2}{S}(C-B)$, $C+\tfrac{a^2+b^2}{S}(A-C)$. The bridge
\decl{Brocard.lemoinePoints_mem} proves that each lies on its side line and on the line through $K$
parallel to another side (the parallel to $BC$ gives points 0 and 1, the parallel to $CA$ points 2
and 3, the parallel to $AB$ points 4 and 5), so they are the classical six points.
\end{enumerate}
\end{definition}

The Brocard weight order is easy to get wrong; the builders checked numerically (30 digits) that
$(1/b^2,1/c^2,1/a^2)$ gives $\angle\Omega AB=\angle\Omega BC=\angle\Omega CA$ and the other cyclic
order gives the mirror point, and Theorem~\ref{thm:brocard} proves it. In Lean, $\angle XYZ$ is
\lean{∠ X Y Z}, the angle at $Y$, with values in $[0,\pi]$; by Mathlib's convention the angle is
$\pi/2$ when $X=Y$ or $Z=Y$.

\section{The Nagel line family}\label{sec:nagel}
\begin{table}[t]
\centering\footnotesize
\caption{Main theorems and their Lean names. Namespaces: \lean{NagelLine} (development
\lean{paperg\_nagel}) and \lean{Brocard} (development \lean{paperg\_brocard}). Starred rows are from the
extension files (\lean{Extensions.lean}). The hypothesis column lists hypotheses beyond the triangle.}
\label{tab:thms}
\begin{tabular}{@{}l p{3.6cm} p{7.4cm} l@{}}
\toprule
Thm & content & Lean & hyp.\\
\midrule
\ref{thm:gergonne} & Gergonne cevians & \decl{gergonnePoint_mem_cevian}, \decl{gergonne_cevians_concurrent} & --\\
\ref{thm:nagel} & Nagel cevians & \decl{nagelPoint_mem_cevian}, \decl{nagel_cevians_concurrent} & --\\
\ref{thm:nagelline} & $N-G=2(G-I)$ & \decl{nagelPoint_vsub_centroid}, \decl{collinear_incenter_centroid_nagelPoint} & --\\
\ref{thm:spieker} & Spieker $=\operatorname{mid}(I,N)$ & \decl{medial_incenter_eq_midpoint}, \decl{collinear_spiekerCenter_incenter_nagelPoint} & --\\
\ref{thm:delong}* & de Longchamps & \decl{midpoint_orthocenter_deLongchampsPoint}, \decl{collinear_deLongchampsPoint_circumcenter_orthocenter}, \decl{collinear_centroid_circumcenter_deLongchampsPoint}, \decl{anticomplementary_orthocenter} & --\\
\ref{thm:medialnagel}* & medial Nagel $=I$ & \decl{medial_nagelPoint} & --\\
\ref{thm:isotomic}* & isotomic contacts & \decl{midpoint_touchpoint_zero}, \lean{\_one}, \lean{\_two} & --\\
\ref{thm:anti}* & $N=$ anticompl.\ incentre & \decl{nagelPoint_eq_anticomplementary_incenter} & --\\
\ref{thm:radical}* & Spieker radical centre & \decl{spiekerCenter_power_exsphere}, \decl{spiekerCenter_power_exsphere_eq}, \decl{spiekerCenter_power_exsphere_eq_inradius} & --\\
\midrule
\ref{thm:symmedian} & symmedians, interior & \decl{symmedian_angles}, \decl{symmedianPoint_mem_interior} & --\\
\ref{thm:pedal} & pedal centroid; converse & \decl{symmedianPoint_eq_centroid_pedal}, \decl{eq_symmedianPoint_of_eq_centroid_pedal} & (ii) $\dim 2$\\
\ref{thm:brocard} & Brocard points & \texttt{brocardPoint\textsubscript{1}\_angles}, \texttt{brocardPoint\textsubscript{2}\_angles} & --\\
\ref{thm:cot} & Brocard angle & \decl{cot_brocardAngle_eq_sum_cot}, \decl{cot_brocardAngle_eq_div_area}, \decl{brocardAngle_pos_le} & --\\
\ref{thm:isogonal}* & isogonal, $|O\Omega|=|O\Omega'|$ & \decl{brocardPoints_isogonal}, \decl{dist_circumcenter_brocardPoint} & --\\
\ref{thm:brocardcircle}* & Brocard circle & \decl{brocardCircle}, \decl{cospherical_brocardCircle}, \decl{brocardCircle_right_angles} & --\\
\ref{thm:degenerate}* & equilateral case & \decl{brocardCircle_degenerate}, \decl{brocard_eq_circumcenter_of_equilateral} & see text\\
\ref{thm:lemoine}* & Lemoine's first circle & \decl{lemoineCircle}, \decl{cospherical_lemoinePoints} & --\\
\ref{thm:droz}* & Droz-Farny & \decl{drozFarny} & see text\\
\bottomrule
\end{tabular}
\end{table}

The theorems of this section are in the development \lean{paperg\_nagel} (library
\lean{NagelLine}); each has the triangle $t$ as its only hypothesis. File:line references are to
the library files at the commit given in Section~\ref{sec:avail}.

\begin{theorem}[Gergonne point]\label{thm:gergonne}
For every triangle and every vertex $i$, the Gergonne point lies on the line through vertex $i$ and
the contact point of the incircle with the opposite side (Mathlib's \lean{t.touchpoint \(\emptyset\) i}).
Consequently the three cevians to the incircle contact points have a common point.
\end{theorem}
Lean: \decl{NagelLine.gergonnePoint_mem_cevian} (\src{paperg_nagel}{NagelLine}{Main}{290}) and
\decl{NagelLine.gergonne_cevians_concurrent} (\src{paperg_nagel}{NagelLine}{Main}{328}).

\begin{theorem}[Nagel point]\label{thm:nagel}
For every triangle and every vertex $i$, the Nagel point lies on the line through vertex $i$ and
the contact point of the excircle opposite vertex $i$ with the side opposite vertex $i$ (Mathlib's
\lean{t.touchpoint \{i\} i}). Consequently the three cevians to the excircle contact points have a
common point. Moreover each such contact point lies strictly between the two other vertices.
\end{theorem}
Lean: \decl{NagelLine.nagelPoint_mem_cevian} (\src{paperg_nagel}{NagelLine}{Main}{334}),
\decl{NagelLine.nagel_cevians_concurrent} (\src{paperg_nagel}{NagelLine}{Main}{357}); the last
sentence is \decl{NagelLine.sbtw_touchpoint_singleton_zero} and its two companions
(\src{paperg_nagel}{NagelLine}{Main}{170}--188).

\begin{theorem}[Nagel line]\label{thm:nagelline}
For every triangle, $N-G=2(G-I)$, where $I$ is Mathlib's incentre. In particular $I$, $G$, $N$ are
collinear.
\end{theorem}
Lean: \decl{NagelLine.nagelPoint_vsub_centroid} (\src{paperg_nagel}{NagelLine}{Main}{362}) and
\decl{NagelLine.collinear_incenter_centroid_nagelPoint} (\src{paperg_nagel}{NagelLine}{Main}{373}).

\begin{theorem}[Spieker centre]\label{thm:spieker}
For every triangle, the incentre of the medial triangle is the midpoint of $I$ and $N$. In
particular the Spieker centre, $I$ and $N$ are collinear.
\end{theorem}
Lean: \decl{NagelLine.medial_incenter_eq_midpoint} (\src{paperg_nagel}{NagelLine}{Main}{385}) and
\decl{NagelLine.collinear_spiekerCenter_incenter_nagelPoint}
(\src{paperg_nagel}{NagelLine}{Main}{408}).

\begin{theorem}[de Longchamps point]\label{thm:delong}
For every triangle, with $L$ the de Longchamps point:
(i) $O$ is the midpoint of $H$ and $L$; (ii) $L$, $O$, $H$ are collinear; (iii) $G$, $O$, $L$ are
collinear; (iv) $L$ is the orthocentre of the anticomplementary triangle.
\end{theorem}
Lean: \decl{NagelLine.midpoint_orthocenter_deLongchampsPoint},
\decl{NagelLine.collinear_deLongchampsPoint_circumcenter_orthocenter},
\decl{NagelLine.collinear_centroid_circumcenter_deLongchampsPoint},
\decl{NagelLine.anticomplementary_orthocenter}
(\src{paperg_nagel}{NagelLine}{Extensions}{228}, 237, 249, 261).
Parts (i) and (ii) are immediate from the definition; (iii) uses Mathlib's Euler line, and (iv) uses
that the circumcentre of the anticomplementary triangle is $H$
(\decl{NagelLine.anticomplementary_circumcenter}).

\begin{theorem}[Nagel point of the medial triangle]\label{thm:medialnagel}
For every triangle, the Nagel point of Mathlib's medial triangle is the incentre $I$ of the triangle.
\end{theorem}
Lean: \decl{NagelLine.medial_nagelPoint} (\src{paperg_nagel}{NagelLine}{Extensions}{274}). (The
statement is written \lean{Triangle.nagelPoint t.medial}, because \lean{Simplex.medial} returns a
\lean{Simplex \(\mathbb{R}\) P 2} and field notation would look for a \lean{Simplex} constant; the term is
the same.)

\begin{theorem}[isotomic contact points]\label{thm:isotomic}
For every triangle and each side, the midpoint of the incircle contact point and the contact point
of the opposite excircle on that side is the midpoint of the side; for the side $BC$,
$\operatorname{mid}(\lean{touchpoint}\ \emptyset\ 0,\ \lean{touchpoint}\ \{0\}\ 0)=\operatorname{mid}(B,C)$,
and likewise for $CA$ and $AB$.
\end{theorem}
Lean: \decl{NagelLine.midpoint_touchpoint_zero}, \decl{NagelLine.midpoint_touchpoint_one},
\decl{NagelLine.midpoint_touchpoint_two} (\src{paperg_nagel}{NagelLine}{Extensions}{297}, 308, 319).

\begin{theorem}[anticomplementary triangle]\label{thm:anti}
For every triangle, the Nagel point is the incentre of the anticomplementary triangle.
\end{theorem}
Lean: \decl{NagelLine.nagelPoint_eq_anticomplementary_incenter}
(\src{paperg_nagel}{NagelLine}{Extensions}{332}).

\begin{theorem}[Spieker centre as radical centre]\label{thm:radical}
For every triangle and all $i,j\in\{0,1,2\}$, the power of the Spieker centre $S_p$ with respect to
the excircle opposite vertex $i$ equals its power with respect to the excircle opposite vertex $j$:
$|S_pJ_i|^2-r_i^2=|S_pJ_j|^2-r_j^2$, where $J_i$ and $r_i$ are the centre and radius of Mathlib's
\lean{t.exsphere \{i\}}. The common value is
\[
\frac14\Bigl(s^2+\frac{(s-a)(s-b)(s-c)}{s}\Bigr)=\frac{s^2+r^2}{4},
\]
with $r$ the inradius.
\end{theorem}
Lean: \decl{NagelLine.spiekerCenter_power_exsphere} (\src{paperg_nagel}{NagelLine}{Extensions}{357})
and, for the value, \decl{NagelLine.spiekerCenter_power_exsphere_eq}
(\src{paperg_nagel}{NagelLine}{Extensions}{201}) for the left-hand form;
\decl{NagelLine.inradius_sq} (Mathlib's \lean{t.inradius} satisfies $r^2=(s-a)(s-b)(s-c)/s$) and
\decl{NagelLine.spiekerCenter_power_exsphere_eq_inradius} for the right-hand form
(\src{paperg_nagel}{NagelLine}{Inradius}{19}, line 40). The last two were added after the Layer-13 comparison
found the right-hand form missing from the library (Section~\ref{sec:verify}).

\section{The symmedian point and the Brocard points}\label{sec:brocard}
The theorems of this section are in the development \lean{paperg\_brocard} (library
\lean{Brocard}). The only hypothesis is the triangle $t$, except in Theorem~\ref{thm:pedal}(ii) and
Theorem~\ref{thm:droz}, where it is stated.

\begin{theorem}[symmedians]\label{thm:symmedian}
For every triangle, with $K$ the symmedian point and $M_a,M_b,M_c$ the midpoints of $BC,CA,AB$:
\[
\begin{aligned}
&\angle BAK=\angle M_aAC,&& \angle CAK=\angle M_aAB,\\
&\angle CBK=\angle M_bBA,&& \angle ABK=\angle M_bBC,\\
&\angle ACK=\angle M_cCB,&& \angle BCK=\angle M_cCA.
\end{aligned}
\]
So at each vertex the line to $K$ is the reflection of the median in the angle bisector, and the
three symmedians pass through $K$. Moreover $K$ is an affine combination of $A,B,C$ with positive
weights, that is, it lies in the relative interior of the triangle (its interior within the plane
of the triangle).
\end{theorem}
Lean: \decl{Brocard.symmedian_angles} (\src{paperg_brocard}{Brocard}{Main}{22}) and
\decl{Brocard.symmedianPoint_mem_interior} (\src{paperg_brocard}{Brocard}{Main}{100}).

\begin{theorem}[pedal triangle]\label{thm:pedal}
(i) For every triangle, $K$ is the centroid of its pedal triangle: if $F_a,F_b,F_c$ are the
orthogonal projections of $K$ onto the lines $BC$, $CA$, $AB$ (Mathlib's
\mdecl{EuclideanGeometry.orthogonalProjection} onto \lean{line[\(\mathbb{R}\), B, C]} etc.), then
$\tfrac13(F_a+F_b+F_c)=K$.
(ii) If moreover $\dim V=2$, then every point $X$ of the plane that is the centroid of its own pedal
triangle is $K$.
\end{theorem}
Lean: \decl{Brocard.symmedianPoint_eq_centroid_pedal} (\src{paperg_brocard}{Brocard}{Main}{115}) and
\decl{Brocard.eq_symmedianPoint_of_eq_centroid_pedal} (\src{paperg_brocard}{Brocard}{Main}{150}),
the latter with the instance hypothesis \lean{[Fact (Module.finrank \(\mathbb{R}\) V = 2)]}.

\begin{theorem}[Brocard points]\label{thm:brocard}
For every triangle, the first Brocard point satisfies
$\angle\Omega AB=\angle\Omega BC=\angle\Omega CA$, and the second satisfies
$\angle\Omega' BA=\angle\Omega' CB=\angle\Omega' AC$.
\end{theorem}
Lean: \texttt{Brocard.brocardPoint\textsubscript{1}\_angles} (\src{paperg_brocard}{Brocard}{Main}{233}) and
\texttt{Brocard.brocardPoint\textsubscript{2}\_angles} (\src{paperg_brocard}{Brocard}{Main}{280}).

\begin{theorem}[Brocard angle]\label{thm:cot}
For every triangle, with $\omega=\angle\Omega AB$ and the vertex angles $\angle BAC$, $\angle ABC$,
$\angle ACB$:
\[
\cot\omega=\cot\angle BAC+\cot\angle ABC+\cot\angle ACB=\frac{a^2+b^2+c^2}{4\Delta},
\qquad 0<\omega\le\frac\pi6.
\]
\end{theorem}
Lean: \decl{Brocard.cot_brocardAngle_eq_sum_cot}, \decl{Brocard.cot_brocardAngle_eq_div_area},
\decl{Brocard.brocardAngle_pos_le} (\src{paperg_brocard}{Brocard}{Main}{327}, 366, 402). Here
$\cot$ is Mathlib's \mdecl{Real.cot}.

\begin{theorem}[isogonal Brocard points]\label{thm:isogonal}
For every triangle, $\angle\Omega AB=\angle\Omega' AC$, $\angle\Omega BC=\angle\Omega' BA$ and
$\angle\Omega CA=\angle\Omega' CB$; and $|O\Omega|=|O\Omega'|$.
\end{theorem}
Lean: \decl{Brocard.brocardPoints_isogonal} (\src{paperg_brocard}{Brocard}{Extensions}{21}) and
\decl{Brocard.dist_circumcenter_brocardPoint} (\src{paperg_brocard}{Brocard}{Extensions}{85}).
``Isogonal'' here means exactly these three angle equalities; no isogonal-conjugation map is defined.
Together with Theorem~\ref{thm:brocard} they say that the two Brocard angles are equal.

\begin{theorem}[Brocard circle]\label{thm:brocardcircle}
For every triangle, each of $O$, $K$, $\Omega$, $\Omega'$ is at distance $\tfrac12|OK|$ from the
midpoint of $OK$; so the four points are cospherical (Mathlib's \mdecl{EuclideanGeometry.Cospherical}).
Moreover $\langle O-\Omega,K-\Omega\rangle=0$ and $\langle O-\Omega',K-\Omega'\rangle=0$, and
$\angle O\Omega K=\angle O\Omega' K=\pi/2$.
\end{theorem}
Lean: \decl{Brocard.brocardCircle}, \decl{Brocard.cospherical_brocardCircle},
\decl{Brocard.brocardCircle_right_angles} (\src{paperg_brocard}{Brocard}{Extensions}{116}, 175, 183).

\begin{theorem}[degenerate Brocard circle]\label{thm:degenerate}
For every triangle, $\Omega=O$ or $\Omega=K$ or $\Omega'=O$ or $\Omega'=K$ or $K=O$ holds if and only
if the triangle is equilateral; and if the triangle is equilateral then $\Omega=\Omega'=K=O$.
\end{theorem}
Lean: \decl{Brocard.brocardCircle_degenerate} (\src{paperg_brocard}{Brocard}{Extensions}{231}) and
\decl{Brocard.brocard_eq_circumcenter_of_equilateral}
(\src{paperg_brocard}{Brocard}{Extensions}{314}), the latter with the two hypotheses $|AB|=|BC|$ and
$|BC|=|CA|$.

\begin{theorem}[Lemoine's first circle]\label{thm:lemoine}
For every triangle, the six Lemoine points (Definition~\ref{def:brocard}) are all at the same distance
from the midpoint of $OK$; in particular they are cospherical.
\end{theorem}
Lean: \decl{Brocard.lemoineCircle} (\src{paperg_brocard}{Brocard}{Extensions}{353}) and
\decl{Brocard.cospherical_lemoinePoints} (\src{paperg_brocard}{Brocard}{Extensions}{439}).

\begin{theorem}[Droz-Farny]\label{thm:droz}
Let $\dim V=2$, let $d_1,d_2\in V$ with $\langle d_1,d_2\rangle=0$, and suppose that neither $d_1$ nor
$d_2$ lies in the direction of any of the three side lines. For $k=1,2,3$ let $X_k$ and $Y_k$ be
points of the $k$-th side line ($BC$, $CA$, $AB$) lying on the line through $H$ with direction $d_1$
and $d_2$ respectively. Then the midpoints of $X_1Y_1$, $X_2Y_2$, $X_3Y_3$ are collinear.
\end{theorem}
Lean: \decl{Brocard.drozFarny} (\src{paperg_brocard}{Brocard}{Extensions}{449}); the lines through $H$
are \lean{AffineSubspace.mk' t.orthocenter (\(\mathbb{R}\) ∙ d\(_i\))}. The non-direction hypotheses
exclude a line through $H$ that is parallel to a side; they also imply $d_1,d_2\neq0$.

\section{Proof method and bridges to Mathlib}\label{sec:method}
\paragraph{Barycentric certificates.} Every point in the statements is, or is proved equal to, an
affine combination of $A,B,C$ with explicit weights that are rational functions of $a,b,c$ (Nagel
family) or of $a^2,b^2,c^2$ (symmedian and Brocard family). Membership of a point in a line, equality
of two points, a midpoint or a collinearity then reduces to an identity between weights, which is
closed by \lean{field\_simp} and \lean{ring}, or by Mathlib's \lean{match\_scalars} where the two sides
are combinations of the same vectors. As in our earlier papers \cite{paperA,paperE}, the identities
and, where needed, their cofactors were found with SymPy and checked numerically (mpmath, 30 digits,
random triangles) before the Lean proofs were written; the numerical checks are evidence, not part of
the proofs.

\paragraph{The Ravi substitution.} For the Nagel family the library proves that there are $x,y,z>0$
with $a=y+z$, $b=z+x$, $c=x+y$ and $s=x+y+z$ (lemma \lean{ravi}, from $s-a,s-b,s-c>0$). After the
substitution the weights $(s-a)/s$, $1/(s-a)$, $a/(2s)$, $(s-b)/a$ and so on become rational functions
of $x,y,z$ with denominators that are visibly positive, so \lean{field\_simp} can clear them without
case analysis. The positivity $s-a>0$ itself is derived from Mathlib's \mdecl{Affine.Simplex.height_pos}
through the squared heights and a factorisation of Heron's polynomial
(\decl{NagelLine.s_sub_a_pos} and its companions).

\paragraph{The touchpoint bridges.} The connection with Mathlib's contact points is made in two steps.
First, Mathlib defines the in- and excentres by weights (\lean{excenterWeights}); we compute them:
the incentre has weights $(a,b,c)/(2s)$ (\decl{NagelLine.incenter_weights}) and
\lean{excenter \{i\}} has weights $\pm a_j/(2(s-a_i))$ with the minus sign at $j=i$
(\decl{NagelLine.excenter_weights}), which identifies it as the excentre opposite vertex $i$. Second,
since \lean{touchpoint} is an orthogonal projection onto a face, a candidate point on the side line is
the contact point exactly when the vector from the centre to it is orthogonal to the side (Mathlib's
\mdecl{EuclideanGeometry.coe_orthogonalProjection_eq_iff_mem}). With the candidate
$B+\frac{s-b}{a}(C-B)$ for the incircle on $BC$, and $B+\frac{s-c}{a}(C-B)$ for the $A$-excircle,
the orthogonality is a polynomial identity in the Gram entries
$|B-A|^2=c^2$, $|C-A|^2=b^2$, $\langle B-A,C-A\rangle=\tfrac12(b^2+c^2-a^2)$. This gives
\decl{NagelLine.touchpoint_empty_zero}, \decl{NagelLine.touchpoint_singleton_zero} and their
companions, after which Theorems~\ref{thm:gergonne}--\ref{thm:isotomic} are weight computations.
The medial triangle's sides are half the original ones (\decl{NagelLine.medial_sideA} etc.), which
transfers the weight formulas to \lean{t.medial}. The anticomplementary triangle is a homothetic image
of $t$, so its centroid, its medial triangle and its circumcentre come from Mathlib's lemmas about
\lean{map} and from \mdecl{EuclideanGeometry.eq_circumcenter_of_dist_eq}; the radical-centre identity
uses Mathlib's formula for the exradius as the inverse of the sum of unnormalised weights.

\paragraph{Angles through inner products.} For the second family all points are written in the
basis $u=B-A$, $v=C-A$ (which need not span $V$; nothing beyond the plane of the triangle is used), so
every inner product is a polynomial in the coordinates and the Gram entries above, whose determinant
is positive (\decl{Brocard.gram_det_pos}). An equality of unoriented angles $\angle XYZ=\angle X'Y'Z'$
is reduced to: the squared cosines are equal, and the two cosine numerators have the same sign
(\decl{Brocard.angle_eq_angle_of_inner}, \decl{Brocard.angle_eq_of_coords}). The first is a rational
identity; for the second the builders found, with SymPy, an expression of the product of the two
numerators as a square times a manifestly positive factor. The cotangent of the angle between vectors $X,Y$ is
$\langle X,Y\rangle/\sqrt{|X|^2|Y|^2-\langle X,Y\rangle^2}$, and for points in coordinates the
square root becomes an absolute $2\times2$ determinant times $\sqrt{|u|^2|v|^2-\langle u,v\rangle^2}$
(\decl{Brocard.cot_angle_eq_inner}, \decl{Brocard.cot_angle_coords}), which turns
Theorem~\ref{thm:cot} into rational identities in $a^2,b^2,c^2$ and $\Delta$. The feet of the
perpendiculars are explicit (\decl{Brocard.orthogonalProjection_vsub}: Mathlib's projection onto
\lean{line[\(\mathbb{R}\), Y, Z]} is the explicit foot); for $K$ the foot on $AB$ is $A+t(B-A)$ with
$t=(-a^2+3b^2+c^2)/(2S)$ with $S=a^2+b^2+c^2$. The converse in Theorem~\ref{thm:pedal} is a linear
system closed with \lean{linear\_combination} and explicit cofactors. For the Brocard circle and the
Lemoine circle the circumcentre and orthocentre are also written in coordinates
(\decl{Brocard.circumcenter_coords}, \decl{Brocard.orthocenter_coords}). In Droz-Farny the
hypothesis $\dim V=2$ is used to write $d_1,d_2$ in the basis $u,v$; the collinearity determinant
factors as a non-vanishing rational factor times a polynomial $Q$, and $Q$ is a multiple of
$\langle d_1,d_2\rangle$.

\section{Degenerate cases and hypotheses}\label{sec:degenerate}
\paragraph{No dimension hypothesis.} Among the numbered theorems, apart from two results, every theorem holds for a triangle in a
real inner product space of any dimension, and carries no \lean{Module.finrank} hypothesis: all the
points involved lie in the plane of the triangle, and the proofs use only coordinates in the basis
$B-A$, $C-A$. The exceptions are Theorem~\ref{thm:pedal}(ii), where $X$ ranges over the whole space
and the proof needs $X$ in the plane of the triangle to write it in coordinates (we have not
formalised the higher-dimensional version); and Theorem~\ref{thm:droz}, where the
directions $d_1,d_2$ must be written in the basis. Both state \lean{[Fact (Module.finrank \(\mathbb{R}\) V = 2)]}
explicitly, and the tampering suite checks that removing it from the pedal converse breaks the build.
Non-degeneracy of the triangle is never a separate hypothesis: it is part of
\mdecl{Affine.Triangle}.

\paragraph{The equilateral Brocard circle.} Theorem~\ref{thm:brocardcircle} has no hypothesis
excluding coincidences. When $\Omega$ coincides with $O$ or $K$, the right-angle statement
$\angle O\Omega K=\pi/2$ holds only because Mathlib defines the angle at a vertex coinciding with an
endpoint to be $\pi/2$; the inner-product form $\langle O-\Omega,K-\Omega\rangle=0$ is then trivially
true. Theorem~\ref{thm:degenerate} says exactly when this happens: for equilateral triangles only, and
then all four points coincide and the ``circle'' has radius $0$. The factor behind both statements
is $a^4+b^4+c^4-a^2b^2-b^2c^2-c^2a^2$, which appears in each of the squared distances. For
non-equilateral triangles the right-angle statement is the classical one, and the Layer-14 witness of
Section~\ref{sec:verify} exhibits a triangle in which $O$, $K$, $\Omega$, $\Omega'$ are pairwise
distinct.

\paragraph{Other edge cases.} The excircle contact points lie on the open sides
(Theorem~\ref{thm:nagel}), so the Nagel cevians are genuine cevians. The Brocard angle bound
$\omega\le\pi/6$ is attained only in the equilateral case; the library proves the non-strict bound,
and a strict instance is checked only on the witness triangle. For right triangles $H$ is a vertex;
the de Longchamps theorems need no exclusion, while in Droz-Farny a line through $H$ could coincide
with a side, which the non-direction hypotheses exclude. Under Lean's convention $x/0=0$ some
expressions would be defined even for degenerate data, but no theorem here relies on it: every
denominator ($s$, $s-a$, $a^2+b^2+c^2$, $a$, $\Delta$, the Gram determinant) is proved positive.

\section{Related work and prior art}\label{sec:related}
\paragraph{Mathlib.} At the pinned revision Mathlib has the Euler line (through the Monge point), the
nine-point circle and its centre on the Euler line, the in- and excentres of a simplex with their
radii, spheres and contact points, the medial simplex, and Ceva's theorem in affine spaces. We found no
Mathlib declaration and no Mathlib pull request mentioning the Nagel, Gergonne, Spieker, de Longchamps,
symmedian (Lemoine) or Brocard points, isogonal conjugation, the Droz-Farny line or Lemoine's circle.
Euler's $OI^2=R(R-2r)$, with excentre versions and $r\le R/2$, is in the open Mathlib pull request
\#44230 \cite{pr44230} (opened 2026-09-26); a proof also appears in the repository \lean{facebookresearch/atlas-lean}. We neither
use nor reprove it.

\paragraph{Earlier Lean work on these centres.} The only related Lean file we found is in the
\lean{lean-genius} repository \cite{leangenius} (file \lean{CevasTheoremOQ04OQ04OQ02.lean}, first committed
2026-06-25, no \lean{sorry}). It defines the Nagel point and the incentre by Ravi-barycentric formulas
in an abstract real module and proves that the defined point lies on the cevian to a foot that is
itself defined by a formula. It is a partial, algebraic form of the cevian statements of
Theorems~\ref{thm:gergonne} and~\ref{thm:nagel}: the feet are not tied to incircle or excircle
tangency and the incentre is not tied to Mathlib's, and there is no Nagel line or ratio statement. We
cite it as partial prior art. We found no Lean formalisation of the other theorems of this paper;
for algebraic Coq forms of some of them see the next paragraph.

\paragraph{Other systems.} In GeoCoq and the Coq high-school geometry library the Euler line is
formalised; we found no Nagel, Gergonne, Spieker, symmedian, Brocard, Lemoine-circle or Droz-Farny
results there, nor in the Archive of Formal Proofs, the HOL~Light distribution or the Mizar
Mathematical Library. Closest to this paper is the certified version of Kimberling's encyclopedia
by Narboux and Braun \cite{certifiedetc}, in Coq. There, triangle centres are defined by their
barycentric or trilinear coordinates, and transformations such as the complement and the isogonal
conjugate by formulas; facts are proved by field normalisation. Its tables list as Coq-verified that
$X(1)$ is the complement of $X(8)$ and $X(10)$ the complement of $X(1)$ (the Nagel line with its
ratio, and the Spieker centre on it), that $X(4)$ is the complement of $X(20)$, and that $X(6)$ is the
isogonal conjugate of $X(2)$. These are algebraic forms of parts of Theorems~\ref{thm:nagelline},
\ref{thm:spieker}, \ref{thm:delong} and~\ref{thm:symmedian}: the points are not tied to incircle or
excircle contact, to an incentre defined by equal distances, or to Euclidean angles, which is what
our statements add. We found no Brocard-point, Brocard-angle, Lemoine-circle or Droz-Farny results
there. All Lean hits for ``brocard'' that we found concern Brocard's number-theoretic
problem $n!+1=m^2$, not the Brocard points.

\paragraph{How the search was done, and its gaps.} The sweep (2026-10-24) grepped a local Mathlib
checkout at the pinned revision, searched mathlib4 pull requests and GitHub Lean code through the
GitHub API, grepped clones of \lean{lean-genius} and \lean{atlas-lean}, ran per-repository searches of
FormalBook, Compfiles and other benchmark collections, and searched local copies of GeoCoq, the Coq
high-school geometry library, HOL~Light and the Mizar library, and the Archive of Formal Proofs
through a GitHub mirror. The query terms were the names of the centres and theorems
(\lean{nagel}, \lean{gergonne}, \lean{symmedian}, \lean{lemoine}, \lean{brocard}, \lean{spieker},
\lean{longchamps}, \lean{isogonal}, \lean{droz}, and others). General web search was unavailable during
that sweep. A second sweep on 2026-10-27 repeated the GitHub searches at the then-current heads of
Mathlib, \lean{atlas-lean}, \lean{formal-conjectures}, Compfiles and \lean{lean-genius}, searched
mathlib4 pull requests opened or merged since 2026-10-01, and added general web search, which found
the Coq work \cite{certifiedetc}. Remaining gaps: ProofAtlas was checked only through its sitemap (752
URLs), not by reading its pages; the certified encyclopedia's website was available only through
archived copies, so its Coq sources and collinearity lists could not be inspected; and GitHub does not
appear to index Mizar files, so the local Mizar search stands. Our statements of novelty are subject
to these gaps.

\paragraph{Classical sources.} The theorems are classical; we follow Johnson \cite{johnson1929},
Coxeter and Greitzer \cite{coxetergreitzer} and Honsberger \cite{honsberger}, and Kimberling's
encyclopedia \cite{etc} for the centre numbers and barycentric coordinates. For the three less
familiar statements: the Spieker centre is the radical centre of the excircles in Yiu
\cite[\S7.3.3, p.~94]{yiu}; the first Lemoine circle and its centre at the midpoint of $KO$ are in
Honsberger \cite[Ch.~9, \S5, pp.~94--95]{honsberger}; the Droz-Farny theorem, stated by Droz-Farny in
1899, is proved synthetically by Ayme \cite{ayme}. Our earlier papers
\cite{paperA,paperB,paperC,paperD,paperE,paperF} use the same certificate method for other
configurations of classical plane geometry.

\section{Verification}\label{sec:verify}
\paragraph{Gates and tampering.} As in our earlier papers, each development has a gate script
(\lean{gate.sh}) that rejects forbidden tokens (\lean{sorry}, \lean{admit}, \lean{native\_decide},
\lean{axiom} declarations and the like), builds the library, and checks with \lean{\#print axioms}
that every declaration in its \lean{REQUIRED} list depends only on \lean{propext},
\lean{Classical.choice} and \lean{Quot.sound}. The \lean{REQUIRED} lists contain every definition,
main theorem and bridge lemma. \lean{paperg\_nagel} passed with 48 declarations (the files
\lean{Basic.lean} and \lean{Main.lean}, Theorems~\ref{thm:gergonne}--\ref{thm:spieker}) and then with
77 after \lean{Extensions.lean} (Theorems~\ref{thm:delong}--\ref{thm:radical}) was added;
\lean{paperg\_brocard} passed with 38 (Theorems~\ref{thm:symmedian}--\ref{thm:cot}) and then 67
after \lean{ExtLemmas.lean} and \lean{Extensions.lean}
(Theorems~\ref{thm:isogonal}--\ref{thm:droz}). The extensions added files and import lines only and
changed no earlier statement. The tampering suite (\lean{tests/tamper.sh}) plants a hidden
\lean{sorry}, an axiom split over two lines, a fake \lean{\#print axioms} report, a
command-redefining macro, a wrapper axiom, and the removal of a needed hypothesis (in
\lean{paperg\_nagel} the triangle binder of \decl{NagelLine.nagelPoint_vsub_centroid} and of
\decl{NagelLine.spiekerCenter_power_exsphere}; in \lean{paperg\_brocard} the dimension hypothesis of
\decl{Brocard.eq_symmedianPoint_of_eq_centroid_pedal} and the orthogonality hypothesis of
\decl{Brocard.drozFarny}); every planted defect was rejected at each stage (7 of 7 in each final
suite). In an early version of \lean{paperg\_nagel} the removed triangle binder was silently
re-supplied by a section \lean{variable}; the tampering suite caught this, and the main theorems were
moved out of that section.

\paragraph{Cross-machine reproduction.} \lean{paperg\_brocard} was built on an arm64 macOS machine from
a copy of the x86-64 Linux repository at the same local commit; the build succeeded, the gate printed
\lean{PASS (67 declarations, standard axioms only)} and the tampering suite rejected all 7 defects. The
SHA-256 hashes of all library \lean{.lean} files of both developments agree between the two machines.
\paragraph{Fresh-clone receipts.} Both repositories were cloned from GitHub at the published commits
and built on both machines (x86-64 Linux and arm64 macOS, Lean v4.33.1, the pinned Mathlib). On each
machine \lean{paperg\_nagel} printed \lean{PASS (151 declarations, standard axioms only)} and
\lean{paperg\_brocard} printed \lean{PASS (147 declarations, standard axioms only)}; the tampering
suites printed \lean{ALL TAMPERING TESTS REJECTED}; and the SHA-256 digest of the sorted library
\lean{.lean} file hashes agreed between the machines (\texttt{201feb97\ldots} and
\texttt{bd4e777b\ldots}). The four receipts are in the supporting material.

\paragraph{Statement-level battery.} The battery of our earlier papers was run on each development
at each stage. The figures are per development; every run's report is in the supporting material
(\lean{reports/battery/}, file names give the machine and the start time).
\begin{itemize}
\item \emph{\lean{paperg\_nagel}, before the extensions} (78 theorems): conclusion mutations 210/210
rejected (7 of them by the heartbeat limit), none surviving; contradiction probe: no contradiction in
12 probed declarations (66 have no hypotheses and were not probed); satisfiability certified for 2 of
78, 76 unresolved; hypothesis ablation: 29 binders needed, 1 unused; linter: no warnings in 2 library
files.
\item \emph{\lean{paperg\_nagel}, with the extensions} (105 theorems): mutations 282/282 rejected (7 by
the heartbeat limit), none surviving; no contradiction in 14 probed declarations (91 without
hypotheses); satisfiability certified for 4 of 105, 101 unresolved; ablation: 32 needed, 1 unused;
linter: no warnings in 3 files (run on the arm64 machine). A run on the x86-64 machine gave the same
mutation, contradiction and satisfiability figures; its ablation stage reported 41 binders needed and
1 unused.
\item \emph{\lean{paperg\_brocard}, before the extensions} (40 theorems): mutations 97/97 rejected, none
by timeout; no contradiction in 13 probed declarations (27 without hypotheses); satisfiability
certified for 8 of 40, 32 unresolved; ablation: 54 needed, 0 unused; linter: no warnings in 2 files.
\item \emph{\lean{paperg\_brocard}, with the extensions} (68 theorems): mutations 152/152 rejected, none
by timeout; no contradiction in 28 probed declarations (40 without hypotheses); satisfiability
certified for 7 of 68, 61 unresolved; ablation: 111 needed, 0 unused; linter: no warnings in 4 files.
The run on the other machine gave the same figures.
\end{itemize}
The single unused binder in \lean{paperg\_nagel} is the hypothesis $0<y+z$ of the auxiliary lemma
\lean{sub\_pos\_aux}, a redundant positivity condition. These figures have limits. The mutation
generator alters only conclusions of certain forms, so a declaration without mutants is not
mutation-tested; most declarations have no hypotheses beyond the triangle and so are not probed for
contradictions; and most satisfiability cases are unresolved by the automatic certifier. For the main
theorems, satisfiability and non-degeneracy are supplied instead by the explicit witnesses of the next
paragraph. Ablation shows which binders the existing proof uses, not which are logically necessary.
With the Layer-14 and Layer-15 files the battery was run again on both machines. For
\lean{paperg\_nagel} at the commit before \lean{Inradius.lean} was added (180 theorems), both runs
rejected 440 of 448 mutations (11 by the heartbeat limit), with 8 surviving and no contradiction;
satisfiability was certified for 30 of 180 on one machine and 32 of 180 on the other (the certifier's
search is time-limited). With \lean{Inradius.lean} (182 theorems) 446 of 454 mutations were rejected,
with the same 8 survivors and no contradiction. For \lean{paperg\_brocard} with the Layer-15 file (147
theorems), two runs on the x86-64 machine and one on the arm64 machine each rejected 285 of 287
mutations, with 2 surviving and no contradiction (one probe was reported invalid by the harness);
satisfiability was certified for 26 to 28 of 147. Besides \lean{sub\_pos\_aux}, the ablation step reports eight unused binders, all in Layer-15 files.
In five classical statements (N3, N4, N6a, N6b, S1b) the hypothesis that the point is on the third
cevian or symmedian is not needed, because the other two already fix the point; the two
nonzero-direction hypotheses of the classical Droz-Farny statement are not used by its proof; and in
the auxiliary lemma \lean{symm\_mem} the hypothesis $0<S$ is redundant. The classical hypotheses are
the classical writer's, kept verbatim. Nine of the ten survivors change a coordinate of a Layer-14
witness triangle (for example $(4,0)\mapsto(5,0)$) inside a non-collinearity lemma; the mutated
triangle is still non-degenerate, so the mutated lemma is true and its survival is expected. The tenth
changes a comment in the Layer-15 Droz-Farny file and has no effect on any statement.
Finally the battery was run on fresh clones of the published commits on both machines (reports
\lean{*\_pub\_*} in the supporting material). \lean{paperg\_nagel} (182 theorems): 446 of 454
mutations rejected (11 by the heartbeat limit), 8 surviving, no contradiction, satisfiability
certified for 30 and 31 of 182; \lean{paperg\_brocard} (147 theorems): 285 of 287 rejected, 2
surviving, no contradiction (one probe invalid), satisfiability certified for 26 and 28 of 147. On
both machines the survivors are the same ten classified above, and the ablation stage reports the same
unused binders.

\paragraph{Statements in Mathlib's language and non-vacuity (Layer~14).} A Layer-14 file in each
development (\lean{NagelLine/L14.lean}, \lean{Brocard/L14.lean}) restates the principal assertions of each main theorem with
Mathlib-only types, the library's points being unfolded into Mathlib terms (explicit affine
combinations with weights written in \lean{dist}, \lean{t.medial.incenter} for the Spieker centre,
$H+2(O-H)$ for the de Longchamps point, and ``any \mdecl{Affine.Triangle} whose vertices are the
images of $A,B,C$ under the homothety with centre $G$ and ratio $-2$'' for the anticomplementary
triangle), and proves it from the library theorem. A check with fully qualified pretty-printing found
no constant of either library in the elaborated types of these 16 + 9 restatements or of the 5 + 5
witness theorems. Each witness theorem proves the hypotheses and the non-degeneracy of the conclusion
on an explicit triangle in \lean{EuclideanSpace \(\mathbb{R}\) (Fin 2)}:
\begin{itemize}
\item For the Nagel family, mainly the 3-4-5 triangle $A=(0,0)$, $B=(4,0)$, $C=(0,3)$: $I=(1,1)$,
$G=(4/3,1)$, $N=(2,1)$, Gergonne point $(8/11,9/11)$, Spieker centre $(3/2,1)$; the cevians are
pairwise distinct lines, $I,G,N$ are pairwise distinct, the incircle and excircle contact points
differ on every side, the anticomplementary triangle has incentre $(2,1)\neq I$, and the excentres
$(6,6)$, $(-2,2)$, $(3,-3)$ with exradii $6,2,3$ give power $37/4$ for the Spieker centre. Since $H=A$ in
this right triangle, the de Longchamps witness uses $(0,0)$, $(4,0)$, $(1,3)$ instead: $O=(2,1)$,
$H=(1,1)$, $L=(3,1)$, pairwise distinct and distinct from $G$.
\item For the symmedian and Brocard family, $A=(0,0)$, $B=(4,0)$, $C=(1,3)$ ($a^2=18$, $b^2=10$,
$c^2=16$, $\Delta=6$): $K=(14/11,12/11)$, $\Omega=(220/157,120/157)$, $\Omega'=(232/157,216/157)$,
$\cot\omega=11/6$ and $0<\omega<\pi/6$; $K$ differs from the vertices and from $G$, and $O,K,\Omega,\Omega'$
are pairwise distinct, so Theorem~\ref{thm:brocardcircle} is non-vacuous; the centroid $G$ is not
the centroid of its own pedal triangle; the six Lemoine points are pairwise distinct; and for
Droz-Farny, with $d_1=(1,2)$, $d_2=(2,-1)$, all hypotheses hold and the three midpoints
$(10/3,2/3)$, $(-2/7,-6/7)$, $(7/4,0)$ are pairwise distinct.
\end{itemize}
Two clauses are not restated there: the strict-betweenness clause of Theorem~\ref{thm:nagel} and the
two forms of the common value in Theorem~\ref{thm:radical}, which are stated only in the library's own
terms. No restatement turned out false or needed a hypothesis that the library theorem lacks. With the
Layer-14 files the gates pass with 91 (\lean{paperg\_nagel}) and 88 (\lean{paperg\_brocard})
declarations on the x86-64 machine, and \lean{paperg\_nagel} also on the arm64 machine; the tampering
suites reject every planted defect.

\paragraph{Blind restatement from this paper (Layer~13).} A version of this paper with every Lean
name, listing, table of declarations and source link removed was given to a fresh AI writer that had
no access to the code; from the prose alone it wrote 36 Lean statements for the 18 numbered theorems,
with its own definitions, compiled with only \lean{sorry} warnings. Its session log shows that it read
only that text and its own files. A comparer then proved, for each pair, the blind statement from the
library statement and conversely, using only unfolding, agreement lemmas between the two sets of
definitions and general Mathlib lemmas, never the library's own theorems. 35 of the 36 pairs are
equivalent in both directions. The exception was Theorem~\ref{thm:radical}: the blind statement
included the inradius form $(s^2+r^2)/4$ written in this paper, which the library did not state, so
the blind statement was strictly stronger. We added the two theorems \decl{NagelLine.inradius_sq} and
\decl{NagelLine.spiekerCenter_power_exsphere_eq_inradius}; with them the blind statement follows by
repackaging. All comparison files compile with no \lean{sorry} and only the standard axioms, and the
library hash was unchanged by the comparison.

\paragraph{Comparison with the classical statements (Layer~15).} A second fresh writer, given only
textbook statements of fifteen classical results (N1--N8: Gergonne, Nagel, Nagel line, Spieker,
isotomic contacts, medial and anticomplementary triangles, de Longchamps, Spieker radical centre;
S1--S7: Lemoine, pedal centroid, Brocard points and angle, isogonality, Brocard circle, first Lemoine
circle, Droz-Farny) and neither this paper nor the code, wrote 23 Lean statements. Comparers then
proved each classical statement from the libraries. All ten Nagel-family statements are covered. The
classical statements are about \emph{any} common point of the cevians, while the library defines the
points by weights and shows they lie on the cevians; the comparer closed this gap with a new
uniqueness proof (two cevians from different vertices meet in at most one point), and also proved the
four-point collinearity of $I$, $G$, $S_p$, $N$, which the library states as two three-point
collinearities. In the symmedian and Brocard family, the pedal-centroid statements and Droz-Farny are
covered directly; the rest were narrower in the library (it proves existence for its explicit points)
and were completed by new proofs: uniqueness of the common point of the symmedians, uniqueness of each
Brocard point in the interior, and the first Lemoine circle for any six points with the defining
incidences. Four classical statements (N4, N6b, N7, N8) carry a non-degeneracy hypothesis on the medial or anticomplementary triangle written by the classical writer; the comparer proved these hypotheses redundant (the medial and anticomplementary triangles of a triangle are non-degenerate). No classical statement disagreed with the library. These proofs, including all new
uniqueness results, were then moved into the libraries as \lean{NagelLine/L15Classical.lean} and
\lean{Brocard/L15Classical.lean}, which work in the plane; with them the gates passed with 149 and
147 declarations, and adding \lean{NagelLine/Inradius.lean} after the Layer-13 comparison raised
\lean{paperg\_nagel} to 151.

\paragraph{Statements checked by hand.} \PH{hand check: the author reads each numbered theorem against
its Lean statement at the published commits}

\section{What is not claimed}\label{sec:notclaimed}
\begin{itemize}
\item We do not claim new mathematics: all theorems are classical.
\item We do not claim Euler's $OI^2=R(R-2r)$, the Euler line, the nine-point circle or any other
result already in Mathlib or in an open Mathlib pull request; we use Mathlib's Euler line in
Theorem~\ref{thm:delong}(iii).
\item We do not claim the first formal treatment of the Nagel and Gergonne cevians in every sense: an
algebraic form with formula-defined feet was formalised earlier in \lean{lean-genius}
(Section~\ref{sec:related}). Nor do we claim the first formal proof of the Nagel line or of the
Spieker centre's position on it in barycentric form, which is among the Coq facts of
\cite{certifiedetc}. Our claim is to the statements with Mathlib's contact points, incentre and
centroid.
\item In the numbered theorems the points are given by their weights, and the theorems prove that
these points have the classical properties (Theorem~\ref{thm:symmedian} by angle equalities at each
vertex); the only uniqueness statement among them is Theorem~\ref{thm:pedal}(ii). The additional
Layer-15 files (Section~\ref{sec:verify}), which work in the plane \lean{EuclideanSpace \(\mathbb{R}\) (Fin 2)},
prove uniqueness of the Nagel point, of the common point of the symmedians and of each interior
Brocard point, and express isogonality by reflected-line incidence. We define no isogonal-conjugation
map.
\item We do not claim novelty beyond the sources searched; see the gaps in Section~\ref{sec:related}.
\end{itemize}

\section{Limitations}\label{sec:limits}
\begin{enumerate}
\item \textbf{Definitions by weights.} The Nagel, Gergonne, symmedian and Brocard points and the
Lemoine points are defined by barycentric formulas. The theorems connect them to Mathlib's contact
points, incentre, centroid, orthocentre, circumcentre and angles, but a reader must check that the
weights are the intended ones; the Layer-14 restatements unfold them so that this check needs only
Mathlib.
\item \textbf{One labelling.} Statements are for the given order $A,B,C$; where a statement has three
cyclic forms, all three are stated (contact points, cevians, symmedian angles), but no general
relabelling lemma is proved.
\item \textbf{Dimension.} The pedal converse and Droz-Farny assume a two-dimensional space. The
additional Layer-15 developments are also planar.
\item \textbf{The equilateral case.} The right-angle form of the Brocard circle is vacuous when the
points coincide, which happens exactly for equilateral triangles (Section~\ref{sec:degenerate}).
\item \textbf{Battery coverage.} Many declarations have no mutants or no hypotheses, and most
satisfiability cases are unresolved by the automatic certifier (Section~\ref{sec:verify}).
\item \textbf{All-AI checking.} The proofs, the numerical checks, the battery and the restatements
were all done by AI agents. \PH{state human checks done before submission (hand check)}
\item \textbf{Prior art.} The search had the gaps listed in Section~\ref{sec:related}.
\end{enumerate}

\section{Availability}\label{sec:avail}
The proofs are in two Lean~4 developments on Lean v4.33.1 with Mathlib pinned at revision
\texttt{0df444a3} (in \lean{lake-manifest.json}). Line counts are of the library \lean{.lean} files,
including the root import file, at the published commits.
\begin{center}\small
\begin{tabular}{lp{8.2cm}rr}
\toprule
repository & library files & files & lines\\
\midrule
\lean{paperg\_nagel} & \lean{Basic} (285), \lean{Main} (414), \lean{Extensions} (362), \lean{Inradius} (45), \lean{L14} (408), \lean{L15Classical} (741) & 7 & 2{,}261\\
\lean{paperg\_brocard} & \lean{Basic} (57), \lean{Lemmas} (338), \lean{Main} (463), \lean{ExtLemmas} (384), \lean{Extensions} (546), \lean{L14} (522), \lean{L15Classical} (723) & 8 & 3{,}040\\
\bottomrule
\end{tabular}
\end{center}
The repositories are \url{https://github.com/jpbald93/paperg_nagel} at commit
\texttt{e40ee54f20588f5181d348767f3f31be20406cc1} and
\url{https://github.com/jpbald93/paperg_brocard} at commit
\texttt{ebb35755a81aec2c7c0673fe93baf8cc5d6eb22d}; each is published as a single commit (the
development history is not included). Every \texttt{File:line} reference in this paper links to
that commit. Supporting material (the prior-art sweep,
the builders' reports, the battery and build reports, the Layer-14 files and reports, and the
fresh-clone receipts) is in the folder \lean{PaperG-triangle-centres} of the repository
\url{https://github.com/jpbald93/LEAN-verification}.
Every name in \lean{\textbackslash decl} font in this paper is in the \lean{REQUIRED} list of the
repository that contains it at the published commit; a script checked all 108 names against the
published \lean{REQUIRED} lists, and all 31 \texttt{File:line} references (and the further line
numbers listed with them) against the declarations at those lines.

\section*{Acknowledgements and use of AI}
Large-language-model AI assistants were used to write the Lean proofs, to run the prior-art search,
the numerical checks and the verification battery, and to draft this paper; parts of the checking were
done by a different model from the one that wrote the proofs. The author directed the work and
reviewed the paper. Every proof is checked by the Lean kernel, and the author takes responsibility for
all content.

\begin{thebibliography}{99}
\bibitem{lean4}
L.~de Moura and S.~Ullrich,
\emph{The Lean 4 theorem prover and programming language},
in: Automated Deduction -- CADE-28, Lecture Notes in Computer Science \textbf{12699}, Springer, 2021, 625--635.

\bibitem{mathlib}
The mathlib Community,
\emph{The Lean mathematical library},
in: Proceedings of the 9th ACM SIGPLAN International Conference on Certified Programs and Proofs (CPP '20), ACM, 2020, 367--381.

\bibitem{pr44230}
W.~Wang, \emph{feat(Geometry/Euclidean): prove Euler's theorem in geometry}, mathlib4 pull request \#44230, opened 26 September 2026 (open at the time of writing).
\url{https://github.com/leanprover-community/mathlib4/pull/44230}

\bibitem{etc}
C.~Kimberling,
\emph{Encyclopedia of Triangle Centers},
\url{https://faculty.evansville.edu/ck6/encyclopedia/ETC.html}.

\bibitem{johnson1929}
R.~A.~Johnson,
\emph{Advanced Euclidean Geometry},
Dover, New York, 1960 (unabridged republication of \emph{Modern Geometry}, Houghton Mifflin, 1929).

\bibitem{coxetergreitzer}
H.~S.~M.~Coxeter and S.~L.~Greitzer,
\emph{Geometry Revisited},
Mathematical Association of America, 1967.

\bibitem{leangenius}
GitHub user \texttt{rjwalters}, \emph{lean-genius} (GitHub repository \url{https://github.com/rjwalters/lean-genius}),
file \texttt{proofs/Proofs/CevasTheoremOQ04OQ04OQ02.lean}, first committed 25 June 2026; consulted at
commit \texttt{470a7df}, 22 October 2026,
\url{https://github.com/rjwalters/lean-genius/blob/470a7df/proofs/Proofs/CevasTheoremOQ04OQ04OQ02.lean}.

\bibitem{yiu}
P.~Yiu, \emph{Introduction to the Geometry of the Triangle}, lecture notes, Florida Atlantic
University, Summer 2001; version 12.1224, December 2012.

\bibitem{ayme}
J.-L.~Ayme, A purely synthetic proof of the Droz-Farny line theorem, \emph{Forum Geometricorum}
\textbf{4} (2004), 219--224.

\bibitem{certifiedetc}
J.~Narboux and D.~Braun, Towards a certified version of the Encyclopedia of Triangle Centers,
\emph{Mathematics in Computer Science} \textbf{10} (2016), 57--73,
\url{https://doi.org/10.1007/s11786-016-0254-4}; website
\url{http://dpt-info.u-strasbg.fr/~narboux/CETC/} (consulted through the Internet Archive, snapshots of
March 2019).

\bibitem{honsberger}
R.~Honsberger,
\emph{Episodes in Nineteenth and Twentieth Century Euclidean Geometry},
New Mathematical Library \textbf{37}, Mathematical Association of America, 1995.

\bibitem{paperA}
J.~Bald,
\emph{Classical plane geometry by coordinate certificates: Desargues, Monge, Simson, Napoleon and Lami in Lean~4},
Zenodo, 2026. \url{https://doi.org/10.5281/zenodo.23239930}

\bibitem{paperB}
J.~Bald,
\emph{The Malfatti circles in Lean~4, with independent certificate proofs of Feuerbach, Miquel and Brahmagupta},
Zenodo, 2026. \url{https://doi.org/10.5281/zenodo.23240500}

\bibitem{paperC}
J.~Bald,
\emph{Barbier's theorem for the Reuleaux triangle and the convex four-vertex theorem in Lean~4, through support functions},
Zenodo, 2026. \url{https://doi.org/10.5281/zenodo.23323062}

\bibitem{paperD}
J.~Bald,
\emph{Circle configurations in classical plane geometry, formalized in Lean~4: Clifford's chain, Miquel's five circles, Lester, Sawayama--Th\'ebault and Kosnita},
Zenodo, 2026. \url{https://doi.org/10.5281/zenodo.23378578}

\bibitem{paperE}
J.~Bald,
\emph{The complete quadrilateral in Lean~4: Newton--Gauss, Bodenmiller, the Miquel point, the Steiner circle and the Simson line},
Zenodo, 2026. \url{https://doi.org/10.5281/zenodo.23407551}

\bibitem{paperF}
J.~Bald,
\emph{The inscribed parabola and the Poncelet point: the complete quadrilateral and the complete quadrangle in Lean~4},
Zenodo, 2026. \url{https://doi.org/10.5281/zenodo.23653615}
\end{thebibliography}
\end{document}
